% 出席確認クイズ 第15回　総括：AI時代の経営情報
%   attendance/attendance15.html から生成（解説は本ガイド用に加筆）。

\quizq{1}{本講義で学んだ「経営情報」の要点として適切なものはどれでしょうか?}%
  {AIにすべて任せれば経営判断は不要}%
  {AIは使わない方がよい}%
  {\correct{情報を経営資源として捉え、AIを検証しながら意思決定に活かす}}%
  {情報は経営とは無関係}%
  {正解は (c)。情報を経営資源として捉え、AIの出力を検証しながら意思決定に活かすことが本講義の中心テーマでした。(a) AI に任せきりにすることも、(b) AI を使わないことも、(d) 情報を経営と無関係とみなすことも、本講義の立場ではありません。}

\quizq{2}{生成AIの出力を確かめる方法として適切なものはどれでしょうか?}%
  {同じAIにもう一度聞くだけ}%
  {長い回答なら信じる}%
  {確かめる必要はない}%
  {\correct{一次情報や公式資料など、別の情報源で裏付けを取る}}%
  {正解は (d)。AIの出力は、一次情報や公式資料など、AIとは独立した情報源で裏付けを取って確かめます。(a) 同じAIに再質問しても同じ誤りを繰り返すことがあり、(b) 回答の長さは正確さの根拠になりません。}

\quizq{3}{AI時代に人が担い続けると考えられる役割として適切なものはどれでしょうか?}%
  {\correct{目的の設定、価値判断、結果への責任}}%
  {計算の代行}%
  {定型作業の反復だけ}%
  {何も担わない}%
  {正解は (a)。目的を定め、価値に基づいて判断し、結果に責任を持つことは、AI ではなく人が担う役割です。(b)(c) は機械に置き換えられやすい作業で、(d) は本講義の立場に反します。}

\quizq{4}{「汎用人工知能(AGI)」の説明として適切なものはどれでしょうか?}%
  {\correct{特定の課題に限らず、幅広い知的作業を人間並みにこなすAI(の概念)}}%
  {表計算ソフトの機能}%
  {画像だけを扱うAI}%
  {すでに一般に普及している家電}%
  {正解は (a)。AGI（Artificial General Intelligence、汎用人工知能）は、特定分野に限らず幅広い知的作業を人間並みにこなせる AI を指す概念で、その定義や実現時期は現在も研究・議論の対象です。(b)(c)(d) は AGI ではありません。}

\quizq{5}{継続的に学ぶ姿勢として適切なものはどれでしょうか?}%
  {流行を無批判に追いかける}%
  {\correct{新しいAIの動向を追いつつ、基本概念と検証の習慣を持ち続ける}}%
  {AIの話題は避ける}%
  {一度学べば更新は不要}%
  {正解は (b)。AIは変化が速いため、動向を追いながらも、基本概念（情報、意思決定、データ、ガバナンス）と出力を検証する習慣を持ち続けることが大切です。(a) 無批判な追随と (d) 更新の放棄はどちらも望ましくなく、(c) 話題を避けることは学びの機会を失います。}
